\begin{theorem}[Levi--Civita 联络的存在唯一性]
给定光滑非退化度量 $g$，存在唯一无挠且满足 $\nabla g=0$ 的仿射联络，由 Koszul 公式 \eqnrefs{eq:math-dg-koszul-formula} 唯一确定。
\end{theorem}

\begin{proof}
要判断度量相容与无挠能否唯一确定联络，可以先用前一条件写出三组导数，再用后一条件消去不需要的联络分量。度量相容给出
\begin{align*}
Xg(Y,Z)&=g(\nabla_XY,Z)+g(Y,\nabla_XZ),\\
Yg(Z,X)&=g(\nabla_YZ,X)+g(Z,\nabla_YX),\\
Zg(X,Y)&=g(\nabla_ZX,Y)+g(X,\nabla_ZY).
\end{align*}
取前两式之和减第三式，得到
\begin{align*}
&Xg(Y,Z)+Yg(Z,X)-Zg(X,Y)\\
={}&g(\nabla_XY,Z)+g(Y,\nabla_XZ)
+g(\nabla_YZ,X)+g(Z,\nabla_YX)\\
&-g(\nabla_ZX,Y)-g(X,\nabla_ZY).
\end{align*}
现在逐次使用无挠关系
\[
\nabla_XZ-\nabla_ZX=[X,Z],
\quad
\nabla_YX-\nabla_XY=[Y,X],
\quad
\nabla_YZ-\nabla_ZY=[Y,Z].
\]
把所有出现的联络项整理到 $g(\nabla_XY,Z)$，得到
\begin{align*}
2g(\nabla_XY,Z)
={}&Xg(Y,Z)+Yg(Z,X)-Zg(X,Y)\\
&+g([X,Y],Z)-g([Y,Z],X)+g([Z,X],Y).
\end{align*}
因此
\begin{equation}\label{eq:math-dg-koszul-formula}
{
\begin{aligned}
2g(\nabla_XY,Z)
={}&Xg(Y,Z)+Yg(Z,X)-Zg(X,Y)\\
&+g([X,Y],Z)-g([Y,Z],X)+g([Z,X],Y).
\end{aligned}}
\end{equation}
右端只含已知的 $g$、向量场与 Lie 括号。由于 $g$ 非退化，\eqnrefs{eq:math-dg-koszul-formula}对每个 $X,Y$ 唯一确定 $\nabla_XY$，故唯一性成立。

反过来，以\eqnrefs{eq:math-dg-koszul-formula}定义 $\nabla_XY$。右端对 $Z$ 是函数线性的，故由非退化性唯一对应一个光滑向量场。
记右端为 $K(X,Y,Z)$。把函数 $a$ 乘入各槽位，利用 $[aX,Y]=a[X,Y]-Y(a)X$ 和 $[X,aY]=a[X,Y]+X(a)Y$，展开后得到
\[
K(aX,Y,Z)=aK(X,Y,Z),\qquad
K(X,aY,Z)=aK(X,Y,Z)+2X(a)g(Y,Z).
\]
所以定义出的导数满足 $\nabla_{aX}Y=a\nabla_XY$ 与 $\nabla_X(aY)=a\nabla_XY+X(a)Y$，确实是仿射联络。
将 $X,Y$ 互换并相减，右端化为 $2g([X,Y],Z)$，故 $T=0$；把 Koszul 公式适当置换后相加，则恢复
\[
Xg(Y,Z)=g(\nabla_XY,Z)+g(Y,\nabla_XZ),
\]
故 $\nabla g=0$。存在性也得到证明。

这就从同一公式证明了存在性和唯一性，且整个构造不依赖坐标或标架。
\end{proof}

\begin{exbox}[用 Koszul 公式求欧氏空间与球面的联络]
将 Koszul 公式用于欧氏空间和半径 $R$ 的球面，比较两种度量的联络系数。

\textbf{解。} 欧氏度量 $g=\delta_{ij}\dd x^i\otimes\dd x^j$ 的分量恒定，坐标基括号为零，故 $\Gamma^k{}_{ij}=0$。
球面度量为 $g=R^2(\dd\theta^2+\sin^2\theta\,\dd\varphi^2)$，唯一非零的度量导数是 $\partial_\theta g_{\varphi\varphi}=2R^2\sin\theta\cos\theta$。
代入坐标联络公式，得到 $\Gamma^\theta{}_{\varphi\varphi}=-\sin\theta\cos\theta$、$\Gamma^\varphi{}_{\theta\varphi}=\Gamma^\varphi{}_{\varphi\theta}=\cot\theta$，其余为零。
常数因子 $R^2$ 被逆度量消去，所以改变球面半径不改变这些坐标联络系数；曲率还需要把曲率二形式与实际面积元比较。
\end{exbox}

\begin{lemma}[截面曲率决定代数曲率张量]
设正定内积空间上的 $R,\widetilde R$ 满足曲率张量的反对称性、配对对称性与第一 Bianchi 恒等式。若每个二维平面的截面曲率相同，则 $R=\widetilde R$。
\end{lemma}

\begin{proof}
设 $R$ 与 $\widetilde R$ 是两个具有 Riemann 曲率全部代数对称性的四阶张量，并且它们给出相同截面曲率。令
\begin{equation*}
A:=R-\widetilde R.
\end{equation*}
则 $A$ 仍具有 Riemann 代数对称性，且对任意单位正交 $u,v$，
\begin{equation*}
A(u,v,v,u)=R(u,v,v,u)-\widetilde R(u,v,v,u)=0.
\end{equation*}
对任意线性无关的 $u,v$，在它们张成的平面内换为正交归一基，四线性与反对称性使 $A(u,v,v,u)$ 等于基上相应数值乘以换基行列式的平方。
对线性相关的 $u,v$，反对称性直接给出零。因此对所有 $u,v$ 都有
\begin{equation*}
A(u,v,v,u)=0.
\end{equation*}
我们证明这已迫使 $A=0$。先把 $v$ 替换为 $v+w$，利用双线性展开和 $A$ 在最后两个槽位的反对称性，
\begin{align*}
0&=A(u,v+w,v+w,u)\\
&=A(u,v,v,u)+A(u,v,w,u)+A(u,w,v,u)+A(u,w,w,u)\\
&=A(u,v,w,u)+A(u,w,v,u),
\end{align*}
即
\begin{equation*}
A(u,v,w,u)+A(u,w,v,u)=0.
\end{equation*}
由配对对称性及两个反对称性，
\[
A(u,w,v,u)=A(v,u,u,w)=-A(u,v,u,w)=A(u,v,w,u).
\]
故上式实际上给出 $2A(u,v,w,u)=0$。
现在把 $u$ 替换为 $u+z$，得到
\[
A(u,v,w,z)+A(z,v,w,u)=0.
\]
又因
\[
A(z,v,w,u)=A(w,u,z,v)=-A(w,u,v,z),
\]
所以 $A(u,v,w,z)=A(w,u,v,z)$，三个循环项相等。
第一 Bianchi 恒等式因而化为
\[
0=A(u,v,w,z)+A(v,w,u,z)+A(w,u,v,z)=3A(u,v,w,z).
\]
我们工作在实数域，故 $A=0$。
这一论证是二次齐次函数 $Q(a,b)=A(a,b,b,a)$ 的两次极化；
$Q$ 分别对 $a,b$ 为二次齐次，而非双线性，不能直接套用双线性型的 $1/4$ 极化公式。

因此在 Riemann 情形，每一点所有二维截面的 $K(\sigma)$ 足以恢复完整的曲率张量；这依赖于反对称性、配对对称性和 Bianchi 恒等式同时成立。
\end{proof}

\begin{proposition}[测地线变分与 Jacobi 方程]
若 $F(s,r)$ 为光滑的仿射参数测地线变分，$T=\partial_sF$、$J=\partial_rF|_{r=0}$，则本书曲率约定下 $\nabla_T^2J+R(J,T)T=0$。
\end{proposition}

\begin{proof}
$F(s,r)$ 是测地线变分：对每个固定 $r$，$s\mapsto F(s,r)$ 都是测地线。沿中心测地线 $\gamma(s)=F(s,0)$ 定义切场 $T=\partial_sF|_{r=0}$ 和变分场 $J=\partial_rF|_{r=0}$。因混合偏导交换，$[T,J]=0$。

由于每条 $r=\mathrm{const.}$ 曲线都是测地线，
\begin{equation*}
 \nabla_TT=0.
\end{equation*}
这是测地线的内蕴定义，等价于 $\dd^2 x^k/\dd s^2+\Gamma^k{}_{ij}\dot x^i\dot x^j=0$。

Levi--Civita 联络无挠，故 $\nabla_TJ-\nabla_JT=[T,J]=0$，即
\begin{equation*}
 \nabla_TJ=\nabla_JT.
\end{equation*}
于是
\begin{equation*}
 \nabla_T\nabla_TJ=\nabla_T\nabla_JT.
\end{equation*}
这一步把"分离场沿测地线的二阶协变导数"转化为"切场沿变分方向的协变导数"，是后续代入曲率定义的关键准备。

按本书曲率约定，
\begin{equation*}
 R(T,J)T
 =\nabla_T\nabla_JT-\nabla_J\nabla_TT-\nabla_{[T,J]}T.
\end{equation*}
后两项分别因 $\nabla_TT=0$（测地线方程）与 $[T,J]=0$（混合偏导交换）消失，故
\begin{equation*}
 R(T,J)T=\nabla_T\nabla_JT=\nabla_T\nabla_TJ.
\end{equation*}

由曲率算子的反对称性 $R(T,J)T=-R(J,T)T$，
\begin{equation*}
 \nabla_T\nabla_TJ=-R(J,T)T,
\end{equation*}
移项即得\eqnrefs{eq:math-dg-jacobi-equation}。
这个方程把变分场的二阶协变导数与曲率作用联系起来，因而可以用曲率判断邻近测地线怎样分离。
\end{proof}

\begin{exbox}[常曲率空间中的横向 Jacobi 场]
沿单位速测地线取平行单位法向场 $e(s)$，求满足 $J(0)=0$、$\nabla_TJ(0)=e(0)$ 的横向 Jacobi 场。

\textbf{解。} 写成 $J(s)=j(s)e(s)$ 后，平行性使协变导数只作用于 $j$。
三种曲率分别给出以下标量方程。
\begin{itemize}
\item 单位球面 $S^n$（截面曲率 $+1$）：方程为 $j''+j=0$，解
\begin{equation*}
 J(s)=\sin(s)\,e(s),
\end{equation*}
$s=\pi$ 处 $J(\pi)=0$，即对径点为共轭点。所有大圆在对径点处汇聚。
\item 双曲空间 $\HH^n$（截面曲率 $-1$）：方程为 $j''-j=0$，解
\begin{equation*}
 J(s)=\sinh(s)\,e(s),
\end{equation*}
$\sinh(s)$ 在 $s>0$ 不为零，无共轭点。测地线指数发散，这是 Cartan--Hadamard 定理的典型实例。
\item $\RR^n$（截面曲率 $0$）：方程为 $j''=0$，解 $J(s)=s\,e(s)$，线性分离，无共轭点。
\end{itemize}
初始分离都为零、初始横向速度相同，正曲率模型在 $s=\pi$ 再次相遇，负曲率模型分离快于平坦模型。这些解给出了曲率比较中的参照函数。
\end{exbox}

Jacobi 方程也适用于伪 Riemann 几何。在广义相对论中，它称为测地偏离方程，描述相邻自由下落粒子的相对加速度；曲率作用 $R(J,T)T$ 对应潮汐效应。引力波改变曲率后，分离矢量 $J$ 的变化可以用于探测。这里的 Lorentz 度量不正定，不能直接把 Riemann 情形的截面曲率聚焦结论逐字套用。

\begin{lemma}[指数映射的 Gauss 引理]
若 $\exp_p$ 在 $v$ 附近有定义，则任意 $w\in T_pM$ 满足\[g_{\exp_p(v)}((d\exp_p)_vv,(d\exp_p)_vw)=g_p(v,w).\]
\end{lemma}

\begin{proof}
要比较指数映射前后的径向内积，可以改变测地线的初速度，并沿测地线跟踪这一内积。取变分
\begin{equation*}
F(t,s)=\exp_p\bigl(t(v+s w)\bigr),
\qquad 0\le t\le1.
\end{equation*}
对固定 $s$，$t\mapsto F(t,s)$ 是测地线（指数映射的定义）。记
\begin{equation*}
T=\partial_tF,
\qquad
J=\partial_sF.
\end{equation*}
$T$ 是测地线切向，$J$ 是变分场。在 $t=0$ 处 $F(0,s)=p$ 不依赖 $s$，故 $J(0)=0$；在 $t=1$ 处 $F(1,s)=\exp_p(v+sw)$，$J(1)=(d\exp_p)_v(w)$。同时 $T(1)=(d\exp_p)_v(v)$。

由于 $[T,J]=0$（混合偏导交换）且 Levi--Civita 联络无挠，$\nabla_TJ=\nabla_JT$。

沿 $t$ 方向求
\begin{equation*}
\frac{\dd}{\dd t}g(T,J)
=g(\nabla_TT,J)+g(T,\nabla_TJ).
\end{equation*}
第一项因测地线方程 $\nabla_TT=0$ 消失。第二项用 $\nabla_TJ=\nabla_JT$：
\begin{equation*}
\frac{\dd}{\dd t}g(T,J)
=g(T,\nabla_JT)
=\frac12 J\bigl[g(T,T)\bigr].
\end{equation*}
最后等式用了度量相容性 $J[g(T,T)]=2g(\nabla_JT,T)$。

每条径向测地线速度平方恒定，并且初速度为 $v+s w$，故
\begin{equation*}
J[g(T,T)]
=\frac{\partial}{\partial s}g_p(v+s w,v+s w)
=2g_p(v+s w,w).
\end{equation*}
关键处用了"测地线速度的范数沿测地线守恒"——这是 $\nabla_TT=0$ 与度量相容性的直接推论：$\frac{\dd}{\dd t}g(T,T)=2g(\nabla_TT,T)=0$。

在 $s=0$ 处积分 $t\in[0,1]$：
\begin{equation*}
g(T(1),J(1))-g(T(0),J(0))
=\int_0^1 g_p(v,w)\,\dd t
=g_p(v,w).
\end{equation*}
又因 $J(0)=0$，初点项 $g(T(0),J(0))=0$，得到
\begin{equation*}
g(T(1),J(1))=g_p(v,w),
\end{equation*}
即 \eqnrefs{eq:math-dg-gauss-lemma}。
取 $w=v$ 时，公式还说明指数映射保持径向长度。这样，曲率比较中的非平凡伸缩都集中在角向方向上。
\end{proof}

\begin{exbox}[Gauss 引理怎样给出极坐标的正交性]
用 Gauss 引理解释欧氏空间与单位球面测地极坐标中的径向、角向正交。

\textbf{解。} 取 $w\perp v$，引理右端为零，所以指数映射把初速度球面的切向量送到与径向方向正交的向量。两种模型具体为：
\begin{itemize}
\item $w\perp v$ 时 $g_p(v,w)=0$，故径向方向 $(d\exp_p)_v v$ 与角向方向 $(d\exp_p)_v w$ 在 $\exp_p(v)$ 处正交。这是"测地极坐标 $g=\dd r^2+r^2 h(r,\theta)$"成立的真正原因。
\item $\RR^n$：$\exp_p(v)=p+v$，$(d\exp_p)_v=\mathrm{id}$，Gauss 引理退化为 $g(v,w)=g(v,w)$，平凡成立。
\item $S^n$：$\exp_p(v)$ 把 $T_pS^n$ 中向量映到球面，Gauss 引理保证球面测地极坐标 $g=\dd r^2+\sin^2(r)\,d\Omega^2$ 中径向与角向正交。
\end{itemize}

因此径向系数固定为 $1$，而球面的 $\sin^2r$ 记录了角向长度相对平坦情形的缩减。
\end{exbox}

\begin{proposition}[测地线能量的二阶变分]
设 $\gamma:[0,L]\to M$ 为仿射参数测地线，$T=\dot\gamma$、$V(0)=V(L)=0$。对 $E=\frac12\int|T|^2\dd s$，\[\delta^2E[V,V]=\int_0^L\bigl(|\nabla_TV|^2-g(R(V,T)T,V)\bigr)\dd s.\]
\end{proposition}

\begin{proof}
对二维变分 $F(s,r)$ 记
\[
T=\partial_sF,
\qquad
V=\partial_rF.
\]
坐标向量场满足 $[T,V]=0$，Levi--Civita 联络无挠，因此
\[
\nabla_VT=\nabla_TV.
\]
对能量求一次导数：
\[
\begin{aligned}
\frac{\dd E}{\dd r}
&=\int_0^L g(\nabla_VT,T)\,\dd s\\
&=\int_0^L g(\nabla_TV,T)\,\dd s\\
&=\bigl[g(V,T)\bigr]_0^L
-\int_0^L g(V,\nabla_TT)\,\dd s.
\end{aligned}
\]
固定端点使边界项消失，得到\eqnrefs{eq:math-dg-energy-first-variation}。

现在假设中心曲线 $r=0$ 是测地线。对一阶变分式关于 $r$ 再求导，并在 $r=0$ 取值。由于 $\nabla_TT=0$，只有 $\nabla_V\nabla_TT$ 的变化留下：
\[
\delta^2E[V,V]
=-\int_0^L g\bigl(V,\nabla_V\nabla_TT\bigr)\,\dd s.
\]
曲率恒等式与 $[V,T]=0$ 给出
\[
\nabla_V\nabla_TT
=\nabla_T\nabla_VT+R(V,T)T
=\nabla_T\nabla_TV+R(V,T)T.
\]
于是
\[
\delta^2E[V,V]
=-\int_0^L g(V,\nabla_T\nabla_TV)\,\dd s
-\int_0^L g(R(V,T)T,V)\,\dd s.
\]
对第一项分部积分，固定端点给出
\[
-\int_0^L g(V,\nabla_T\nabla_TV)\,\dd s
=\int_0^L|\nabla_TV|^2\,\dd s.
\]
故
\begin{equation}
{
\delta^2E[V,V]
=\int_0^L
\left(
|\nabla_TV|^2-g(R(V,T)T,V)
\right)\dd s.}
\label{eq:math-dg-energy-second-variation}
\end{equation}
极化后即得到\eqnrefs{eq:math-dg-index-form}。
\end{proof}

指标形式 $I(V,V)=\int(|\nabla_TV|^2-g(R(V,T)T,V))\,\dd s$ 的两项竞争：协变导数项 $|\nabla_TV|^2\ge0$ 使测地线稳定，曲率项 $-g(R(V,T)T,V)$ 在正曲率方向为负使测地线不稳定。

若存在 $V\neq0$ 使 $I(V,V)<0$，则测地线不是能量极小——存在更短的邻近曲线。

若有非零 Jacobi 场 $J$ 且 $J(0)=J(L)=0$，则终点与起点共轭，且 $I(J,J)=0$。

反过来，仅有 $I(V,V)=0$ 不能断言 $V$ 为 Jacobi 场，须要求 $I(V,W)=0$ 对所有固定端点场 $W$ 成立；越过首个共轭点后出现负方向，最短性失效。

\begin{exbox}[球面测地线越过对跖点后为何失去最短性]
在 $\mathbb S^n(R)$ 上取单位速大圆测地线 $\gamma:[0,L]\to\mathbb S^n(R)$，用指标形式判断长度为 $L$ 的曲段何时出现负方向。

\textbf{解。} 沿 $\gamma$ 取法向变分场 $V=f(s)\vb e_\perp$，其中 $\vb e_\perp$ 是平行单位法向量。球面截面曲率 $K=1/R^2$，故
\[
 I(V,V)=\int_0^L\!\left(|f'|^2-\frac{1}{R^2}|f|^2\right)\dd s.
\]
取 $f(s)=\sin(\pi s/L)$，它对每个 $L$ 都满足固定端点条件。利用正弦、余弦平方在 $[0,L]$ 上的积分均为 $L/2$，得到
\[
I(V,V)=\frac L2\left(\frac{\pi^2}{L^2}-\frac1{R^2}\right).
\]
当 $L=\pi R$ 时，$V$ 是非零 Jacobi 场，且 $I(V,V)=0$；终点为对跖点，仍有实现距离的大圆半弧，但最短曲线不唯一。
当 $L>\pi R$ 时，同一个测试场给 $I(V,V)<0$，所以曲段已不再局部最短。
\end{exbox}

\begin{theorem}[Bonnet--Myers 直径界]
设 $(M^n,g)$ 连通、完备，$n\ge2$，$\operatorname{Ric}\ge(n-1)kg$、$k>0$。则 $\operatorname{diam}M\le\pi/\sqrt k$，且 $M$ 紧致。
\end{theorem}

\begin{proof}
若两点相距过远，沿它们的最短测地线可以构造一个使能量下降的变分，从而与最短性矛盾。反设两点距离为
\[
L>\frac{\pi}{\sqrt k}.
\]
由 Hopf--Rinow，存在连接它们的单位速最短测地线
\[
\gamma:[0,L]\to M.
\]
沿 $\gamma$ 取平行正交标架
\[
E_1,\ldots,E_{n-1},E_n=T:=\dot\gamma.
\]
对 $i=1,\ldots,n-1$ 定义固定端点变分场
\[
V_i(t)
=\sin\!\left(\frac{\pi t}{L}\right)E_i(t).
\]
由于整段 $\gamma$ 实现两端距离，它在全部固定端点变分下使能量局部极小，指标形式必须非负：
\[
I(V_i,V_i)\ge0.
\]
直接计算
\[
|\nabla_TV_i|^2
=\frac{\pi^2}{L^2}
\cos^2\!\left(\frac{\pi t}{L}\right),
\]
以及
\[
g(R(V_i,T)T,V_i)
=\sin^2\!\left(\frac{\pi t}{L}\right)
K(T,E_i).
\]
对 $i$ 求和，并使用
\[
\sum_{i=1}^{n-1}K(T,E_i)
=\operatorname{Ric}(T,T)
\ge(n-1)k,
\]
得到
\[
\begin{aligned}
0
&\le\sum_{i=1}^{n-1}I(V_i,V_i)\\
&\le(n-1)
\int_0^L
\left[
\frac{\pi^2}{L^2}
\cos^2\!\left(\frac{\pi t}{L}\right)
-k\sin^2\!\left(\frac{\pi t}{L}\right)
\right]\dd t.
\end{aligned}
\]
利用
\[
\int_0^L\sin^2\!\left(\frac{\pi t}{L}\right)\dd t
=
\int_0^L\cos^2\!\left(\frac{\pi t}{L}\right)\dd t
=\frac L2,
\]
可得
\[
0
\le
\frac{(n-1)L}{2}
\left(
\frac{\pi^2}{L^2}-k
\right).
\]
但 $L>\pi/\sqrt k$ 使括号为负，矛盾。因此任何两点距离都不超过 $\pi/\sqrt k$，得到\eqnrefs{eq:math-dg-bonnet-myers-diameter}。
固定基点 $p$，整个流形就是闭球 $\overline B_{\pi/\sqrt k}(p)$；由完备性和 Hopf--Rinow，这个闭有界集合紧，故 $M$ 紧致。
这里使用的是统一的正 Ricci 下界，而没有对每个截面分别要求正下界。
\end{proof}

\begin{exbox}[Bonnet--Myers 的等号模型与适用范围]
比较球面、平坦环面与双曲空间能否满足定理的曲率假设。

\textbf{解。} 半径 $R$ 的球面 $\mathbb S^n(R)$ 满足 $\operatorname{Ric}=(n-1)R^{-2}g$。取 $k=R^{-2}$ 得直径 $\le\pi R$；一对对跖点的距离为 $\pi R$，故等号达到。

平坦环面 $T^n$ 满足 $\operatorname{Ric}=0$，不存在定理要求的 $k>0$。它仍是紧致的，但可以通过整体缩放把直径取得任意大，故零下界不能给出统一的直径上界。

单位曲率尺度的双曲空间满足 $\operatorname{Ric}=-(n-1)g$，也不满足正下界假设，且直径无限。这两个模型说明，定理假设失效时不能保留其统一直径界。
\end{exbox}

\begin{theorem}[Bishop--Gromov 体积单调性]
设 $(M^n,g)$ 完备，$\operatorname{Ric}\ge(n-1)kg$。对任意 $p$，$\operatorname{Vol}B_p(r)/V_k(r)$ 随 $r$ 单调不增；$k\le0$ 时取 $r>0$，$k>0$ 时取 $0<r<\pi/\sqrt k$。
\end{theorem}

\begin{proof}
先比较每个径向方向上的体积伸缩，再对方向和半径积分，才能得到球体积比的单调性。在指数映射的极坐标中，体积元可写成
\[
\operatorname{vol}_g=J(r,\theta)\,\dd r\,\dd\theta,
\]
其中 $J(r,\theta)$ 是 $\dd\exp_p$ 在横向方向的 Jacobi 行列式。由行列式微分公式，
\[
\partial_r\log J(r,\theta)
=\tr A_r
=\Delta_{\mathrm{LB}} r.
\]
结合\eqnrefs{eq:math-dg-laplacian-comparison}，
\[
\partial_r
\log\frac{J(r,\theta)}{S_k(r)^{n-1}}
\le0.
\]
又因 $r\downarrow0$ 时
\[
J(r,\theta)\sim r^{n-1},
\qquad
S_k(r)^{n-1}\sim r^{n-1},
\]
所以对每个未越过割时刻的方向，
\begin{equation}
\frac{J(r,\theta)}{S_k(r)^{n-1}}
\quad\text{随 $r$ 单调不增。}
\label{eq:math-dg-jacobian-comparison}
\end{equation}
为处理割点，把每条方向在割时刻以后的 Jacobi 密度延拓为零。
割迹集合测度为零，测地极坐标积分因此仍给出球体积。
记单位方向球面积为 $a_{n-1}$，则
\[
q(r)=\frac1{a_{n-1}}\int_{S_pM}
\frac{J(r,\theta)\mathbf1_{r<c(\theta)}}{S_k(r)^{n-1}}\dd\theta
\]
单调不增；$c(\theta)$ 是该方向的割时刻。
球体积比正是其加权平均：
\[
\frac{\operatorname{Vol}B_p(r)}{V_k(r)}
=\frac{\int_0^r q(t)S_k(t)^{n-1}\dd t}
{\int_0^r S_k(t)^{n-1}\dd t}.
\]
记这个加权平均为 $Q(r)$，且令 $D(r)=\int_0^rS_k(t)^{n-1}\,\dd t$。在几乎处处可微的半径上，商法则给
\[
Q'(r)=\frac{S_k(r)^{n-1}}{D(r)}\bigl(q(r)-Q(r)\bigr)\le0.
\]
不等号来自 $q$ 单调不增：末端值不超过此前的加权平均。$Q$ 局部绝对连续，故这个导数不等式给出全部允许半径上的单调性。
这样才从逐方向 Jacobi 比较得到球体积比的单调性，而不只是体积上界。若 $V_k(r)$ 表示常截面曲率 $k$ 的单连通模型空间中半径 $r$ 的球体积，则
\begin{equation}
{
\frac{\operatorname{Vol}_g(B_p(r))}{V_k(r)}
\quad\text{随 $r$ 单调不增。}
}
\label{eq:math-dg-bishop-gromov}
\end{equation}
特别地，当 $\operatorname{Ric}\ge0$ 时，
\[
\frac{\operatorname{Vol}_g(B_p(r))}{\omega_n r^n}
\]
单调不增且不超过 $1$。因此正或非负 Ricci 曲率不仅限制测地线聚焦，也限制体积增长速度。
\end{proof}
